% Part: normal-modal-logic
% Chapter: tableaux
% Section: simple-S5

\documentclass[../../../include/open-logic-section]{subfiles}

\begin{document}

\olfileid{nml}{tab}{s5}

\olsection{\Log{S5} માટે સરળ \usetoken{P}{tableau}}

\Log{S5} સાર્વત્રિક નિદર્શનોના વર્ગને અનુલક્ષીને યથાર્થ અને પૂર્ણ
છે, એટલે એવા નિદર્શનો જેમાં દરેક જગત દરેક જગતમાંથી સુલભ
હોય છે. સાર્વત્રિક નિદર્શનમાં સુલભતા-સંબંધથી કોઈ વધારાનો
ભેદ પડતો નથી: ``એવું જગત~$w$ છે જ્યાં
$\mSat{M}{!A}[w]$'' ત્યારે અને માત્ર ત્યારે સાચું છે જ્યારે
એવું જગત~$w$~$u$થી સુલભ હોય. તેથી \Log{S5}માં નિદર્શનને
જગતોના ગણ અને નિમણૂક~$V$ વડે જ વ્યાખ્યાયિત કરી શકીએ.
આથી !!{tableau}ના નિયમો પણ સરળ કરી શકાય એવું સૂચન મળે છે.
સામાન્ય કિસ્સામાં ઉપસર્ગ તરીકે ધન પૂર્ણાંકોની શ્રેણીઓ લઈએ
છીએ, જેથી કયા ઉપસર્ગ બીજા કયા ઉપસર્ગથી સુલભ જગતોને નામ આપે
છે તેનો હિસાબ રાખી શકાય: $\sigma.n$ એ~$\sigma$થી સુલભ
જગતનું નામ આપે છે. પરંતુ \Log{S5}માં દરેક જગત બીજા દરેક
જગતથી સુલભ છે, તેથી આવો હિસાબ રાખવો જરૂરી નથી. તેના
બદલે ધન પૂર્ણાંકોને જ ઉપસર્ગ તરીકે વાપરી શકીએ. સરળ બનાવેલા
નિયમો \olref{tab:rules-S5}માં આપ્યા છે.

\begin{table}
  \begin{center}
    \def\arraystretch{3}\def\fCenter{}
    \begin{tabular}{|c|c|}
    \hline
    \iftag{prvBox}{
      \AxiomC{\sFmla{\True}{\Box !A}[n]}
      \RightLabel{\TRule{\True}{\Box}}
      \UnaryInfC{\sFmla{\True}{!A}[m]}
      \DisplayProof
      &
      \AxiomC{\sFmla{\False}{\Box !A}[n]}
      \RightLabel{\TRule{\False}{\Box}}
      \UnaryInfC{\sFmla{\False}{!A}[m]}
      \DisplayProof\\
      $m$ વપરાયેલો છે & $m$ નવો છે
      \\[1ex]
      \hline}{}
    \iftag{prvDiamond}{
      \AxiomC{\sFmla{\True}{\Diamond !A}[n]}
      \RightLabel{\TRule{\True}{\Diamond}}
      \UnaryInfC{\sFmla{\True}{!A}[m]}
      \DisplayProof
      &
      \AxiomC{\sFmla{\False}{\Diamond !A}[n]}
      \RightLabel{\TRule{\False}{\Diamond}}
      \UnaryInfC{\sFmla{\False}{!A}[m]}
      \DisplayProof\\
      $m$ નવો છે & $m$ વપરાયેલો છે
      \\[1ex]
      \hline}{}
    \end{tabular}
  \end{center}
  \caption{\Log{S5} માટેના સરળ બનાવેલા નિયમો.}
  \ollabel{tab:rules-S5}
\end{table}

\begin{ex}
  $\Log{S5} \Proves \Ax{5}$, એટલે કે,
  $\Diamond!A \lif \Box\Diamond!A$ બતાવતું સરળ બનાવેલું બંધ
  ટેબ્લો આપીએ છીએ.
  \begin{oltableau}
    [\pFmla{\False}{\Diamond\formula{A} \lif \Box\Diamond \formula{A}}{1},
      just = \TAss
      [\pFmla{\True}{\Diamond \formula{A}}{1}, just = {\TRule{\False}{\lif}[1]}
        [\pFmla{\False}{\Box\Diamond \formula{A}}{1},
          just = {\TRule{\False}{\lif}[1]}
          [\pFmla{\False}{\Diamond \formula{A}}{2},
            just = {\TRule{\False}{\Box}[3]}
            [\pFmla{\True}{\formula{A}}{3},
              just = {\TRule{\True}{\Diamond}[2]}
                [\pFmla{\False}{\formula{A}}{3},
                  just = {\TRule{\False}{\Diamond}[4]}, close]
            ]
          ]
        ]
      ]
    ]
  \end{oltableau}
\end{ex}

\end{document}
